\documentclass[12pt,a4paper]{report}

% Font and encoding
\usepackage[T2A, T1]{fontenc}
\usepackage[utf8]{inputenc}
\usepackage[russian, english]{babel}
\babelprovide[import]{kazakh}
\usepackage{tempora}
\usepackage{geometry}
\geometry{a4paper, margin=2.5cm}

% Packages
\usepackage{setspace}
\usepackage{graphicx}
\usepackage{tikz}
\usetikzlibrary{arrows.meta,positioning,shapes.geometric,calc}
\usepackage{array}
\usepackage{tabularx}
\usepackage{booktabs}
\usepackage[table]{xcolor}
\usepackage{hyperref}
\usepackage{listings}
\usepackage{caption}
\usepackage{titlesec}
\usepackage{tocloft}
\usepackage{indentfirst}
\usepackage{microtype}
\usepackage{parskip}

% Spacing
\onehalfspacing

\hypersetup{
    colorlinks=true,
    linkcolor=black,
    urlcolor=blue,
    citecolor=black
}

\lstset{
    basicstyle=\ttfamily\small,
    breaklines=true,
    frame=single,
    backgroundcolor=\color{gray!10},
    keywordstyle=\color{blue},
    stringstyle=\color{red!70!black},
    commentstyle=\color{green!50!black},
    showstringspaces=false,
    numbers=left,
    numberstyle=\tiny,
    tabsize=2
}

\titleformat{\chapter}[display]
  {\normalfont\Large\bfseries}{\chaptertitlename\ \thechapter}{20pt}{\LARGE}
\titlespacing*{\chapter}{0pt}{20pt}{20pt}

\titleformat{\section}
  {\normalfont\large\bfseries}{\thesection}{1em}{}
\titleformat{\subsection}
  {\normalfont\normalsize\bfseries}{\thesubsection}{1em}{}

\definecolor{tableblue}{RGB}{129, 183, 255}

% -----------------------------------------------
\begin{document}

% -----------------------------------------------
% TITLE PAGE 1
% -----------------------------------------------
\begin{titlepage}
\centering

\includegraphics[width=0.3\textwidth]{sdu_logo.png}\\[1cm]

{\large SDU University\\Faculty of Engineering and Natural Sciences}

\vspace{3cm}
{\Huge\bfseries DIPLOMA PROJECT}

\vspace{0.5cm}
{\Large ARIA. The personal AI assistant for home monitoring, device control, and information access.}

\vspace{1.5cm}
{\normalsize
Kenesbayev Arsen, Orynbek Aidos, Sagibek Adil\\[0.5em]
6B06102   -- ``Computer Science''
}

\vfill
{\normalsize Kaskelen, 2026}
\end{titlepage}

% -----------------------------------------------
% TITLE PAGE 2
% -----------------------------------------------
\begin{titlepage}
    \centering

    \includegraphics[width=0.25\textwidth]{sdu_logo.png}\\[0.5cm]

    {\Large \textbf{SDU University}}\\[0.2cm]
    {\large Faculty of Engineering and Natural Sciences}

    \vfill

    \begin{flushright}
        \begin{minipage}{0.4\textwidth}
            \raggedleft
            Dean of Faculty\\
            Assistant Professor,\\
            Ph.D. Ramis Akhmedov\\[1.5cm]

            \rule{7cm}{0.6pt}\\
            \guillemotleft\hspace{1cm}\guillemotright\hspace{2cm} 2026y.
        \end{minipage}
    \end{flushright}

\vspace{2cm}
\centering
{\Huge\bfseries DIPLOMA PROJECT}

\vspace{0.5cm}
{\Large ARIA. The personal AI assistant for home monitoring, device control, and information access.}

\vspace{1cm}
{\normalsize 6B06102  -- ``Computer Science''}

\vspace{1.5cm}
\begin{tabular}{ll}
\textbf{Supervisor:}    & Binara Imankulova \\[0.5em]
\textbf{Students:}      & Kenesbayev Arsen, Orynbek Aidos, \\
                        & Sagibek Adil \\
\end{tabular}

\vfill
{\normalsize Kaskelen, 2026}
\end{titlepage}

% -----------------------------------------------
% ABSTRACT (English)
% -----------------------------------------------
\chapter*{Abstract}
\addcontentsline{toc}{chapter}{Abstract}

This thesis describes \textbf{ARIA}, an integrated smart-home assistant that combines a locally hosted web application with a custom 3D-printed device. A Flask server provides a browser dashboard, conversational AI (Google Gemini), optional Kazakh-oriented tooling, Gmail integration via OAuth~2.0, weather information, YouTube-based music playback, and control of a Yeelight smart bulb. Physical interaction uses two Espressif ESP32 modules on the same Wi-Fi network: an ESP32-CAM delivers a live MJPEG video stream and drives a stepper motor for horizontal panning (a tilt servo is present in firmware but was not reliably usable in the printed mechanical assembly); a second ESP32 module bridges a digital I2S microphone (INMP441) and a class-D amplifier (MAX98357) to the PC using low-latency UDP audio. Speech recognition uses \textit{faster-whisper}; synthesis uses \textit{edge-tts}; an optional wake-word model (\textit{openWakeWord}) can trigger hands-free sessions. The work highlights how commodity microcontrollers, open Python libraries, and cloud LLM APIs can be composed into a single, extensible home system while keeping part of the stack under the user's control.

% -----------------------------------------------
% ABSTRACT (Kazakh)
% -----------------------------------------------
\begin{otherlanguage}{russian}
\chapter*{Аңдатпа}
\addcontentsline{toc}{chapter}{Аңдатпа}
Бұл дипломдық жобада үйді бақылау, құрылғыларды басқару және ақпаратқа қол жеткізу мақсатында жасалған \textbf{ARIA} интеграцияланған ақылды үй көмекшісі сипатталады. Flask веб-қосымшасы браузер арқылы басқару тақтасын, сұхбаттасу үшін Google Gemini LLM, OAuth~2.0 арқылы Gmail, ауа райы, YouTube негізіндегі музыка және Yeelight шамын басқаруды ұсынады. Физикалық деңгейде екі ESP32 модульі бір Wi-Fi желісінде жұмыс істейді: ESP32-CAM тікелей бейне ағынын және горизонтальды панорамалау үшін шагты қозғалтқышты басқарады (бұрау сервомеханизмі кодта бар, бірақ 3D басып шығару механикасында толық пайдаланылмаған); екінші ESP32 INMP441 микрофоны мен MAX98357 динамигін UDP арқылы дербес компьютермен байланыстырады. Сөйлеуді мәтінге айналдыру \textit{faster-whisper}, синтез \textit{edge-tts}, қосымша ояту сөзі \textit{openWakeWord} арқылы іске асырылады.
\end{otherlanguage}

% -----------------------------------------------
% ABSTRACT (Russian)
% -----------------------------------------------
\begin{otherlanguage}{russian}
\chapter*{Аннотация}
\addcontentsline{toc}{chapter}{Аннотация}
В данной работе описывается интегрированный помощник для умного дома \textbf{ARIA}, объединяющий локально развернутое веб-приложение и аппаратный модуль с 3D-печатным корпусом. Сервер на Flask предоставляет панель управления в браузере, диалог на базе большой языковой модели Google Gemini, интеграцию с Gmail по OAuth~2.0, данные о погоде, воспроизведение музыки через YouTube и управление умной лампой Yeelight. На аппаратном уровне используются два модуля Espressif ESP32 в одной сети Wi-Fi: ESP32-CAM обеспечивает потоковое MJPEG-видео и управляет шаговым двигателем для горизонтального панорамирования (сервопривод наклона реализован в прошивке, но в напечатанной механике не был устойчиво задействован); второй ESP32 связывает цифровой микрофон INMP441 и усилитель MAX98357 с персональным компьютером по UDP. Распознавание речи выполняется библиотекой \textit{faster-whisper}, синтез речи --- \textit{edge-tts}, опционально используется детектор ключевого слова \textit{openWakeWord}.
\end{otherlanguage}

\tableofcontents
\listoffigures
\listoftables

% -----------------------------------------------
\chapter{Introduction}

\section{Problem background}

Consumer demand for connected homes continues to grow: lighting, security cameras, voice assistants, and media devices are increasingly networked. Analyst summaries of the global smart-home segment project sustained revenue expansion over multi-year horizons; as one illustrative industry estimate, Technavio reported expected incremental growth on the order of USD~255.2~billion for the worldwide market between 2025 and 2029, with strong CAGR figures cited for the same interval \cite{technavio}. Independent market studies quote different absolute market sizes depending on segmentation (hardware only versus services), but the directional trend is consistent: households adopt more IP-connected appliances and expect them to interoperate through a phone or voice interface.

Parallel to this adoption, users face recurring tensions. \textbf{Closed ecosystems} tie automations to a single vendor account and make it difficult to combine, for example, a generic IP camera, a third-party bulb, and a desktop LLM workflow. \textbf{Privacy and data residency} matter for email and indoor video: many users prefer a model where credentials and recent message text stay on a machine they control, while only anonymized prompts hit a cloud API. \textbf{Multilingual support} remains uneven; assistants tuned primarily for English and a few majority languages under-serve Kazakh--Russian bilingual households common in Central Asia. Finally, \textbf{integration cost} is high for hobbyists: stitching HTTP camera streams, UDP audio, OAuth, and semantic intent parsing into one coherent interface requires full-stack skills that off-the-shelf gadgets do not teach.

Commercial smart speakers and cloud assistants offer polished experiences but often lock integrations to vendor APIs and require always-on cloud policies. Educational and maker-oriented robots emphasize mechanics or curriculum, not necessarily a full web stack with email, weather, and third-party bulbs. Academic and industrial IoT literature frequently recommends \textit{edge--cloud partitioning}: keep sensing and actuation at the edge, run heavy inference or orchestration on a gateway \cite{esp-idf}. ARIA follows that pattern with two ESP32 edge nodes (vision/motion and audio) and a PC-hosted Flask application as the gateway.

The project addresses the gap by implementing an \textbf{open, modular} assistant where the user hosts the control plane on a PC, keeps structured data (chat and cached mail metadata) in a local SQLite database, and attaches inexpensive ESP32 peripherals for vision and audio I/O. The physical device reinforces the product narrative: a recognizable ``head'' with a camera aperture and status lighting, 3D-printed as an iterative design artifact rather than an anonymous bench of breakout boards.

\section{Scope and assumptions}

The thesis focuses on a \textbf{working prototype} on a single LAN: one Flask host, one Yeelight bulb (model YLDP005, Table~\ref{tab:yeelight-bulb}), one ESP32-CAM module, and one ESP32-WROOM audio bridge. Multi-user cloud deployment, mobile apps outside the browser, and hardened TLS are explicitly out of scope but noted as future work. The assistant relies on external APIs (Gemini, Gmail, YouTube, OpenWeather) where keys are available; behaviour degrades gracefully when a service is disabled.

Table~\ref{tab:requirements} summarizes major functional expectations that guided acceptance testing during development.

\begin{table}[htbp]
\centering
\caption{High-level functional requirements for the ARIA prototype.}
\label{tab:requirements}
\small
\begin{tabular}{@{}p{0.22\textwidth}p{0.68\textwidth}@{}}
\toprule
\textbf{Area} & \textbf{Expected behaviour} \\
\midrule
Dashboard UI & Web UI with sections for dashboard, camera, functions, settings; status indicators for server online. \\
Conversation & User messages answered via Gemini; history persisted per chat session in SQLite. \\
Voice & Optional capture from ESP32 I2S mic, STT with faster-whisper, TTS with edge-tts; optional wake word. \\
Vision & MJPEG stream from ESP32-CAM embedded in UI; pan/tilt commands from UI (tilt may be limited mechanically). \\
Lighting & Yeelight on/off, brightness, named colors, warm/cool white, RGB cycle from UI and parsed chat phrases. \\
Information & Weather and forecast cards; Gmail sync and listing when OAuth configured; YouTube music search. \\
\bottomrule
\end{tabular}
\end{table}

\section{Aim and objectives}

\textbf{Aim.} To design and implement ARIA: a multimodal smart-home assistant that integrates a Flask-based application, cloud LLM reasoning, local speech processing, smart-light control, and an ESP32-based sensing/actuation head fed from a single 5~V supply.

\textbf{Objectives.}
\begin{itemize}
    \item Build a web dashboard for chat, settings, camera viewing, and auxiliary functions, with persistent chat history and optional Gmail linking.
    \item Integrate Google Gemini for natural-language dialogue, with configurable assistant personality and language detection hooks.
    \item Implement STT (faster-whisper) and TTS (edge-tts) on the host, streaming audio to and from a dedicated ESP32 audio bridge (INMP441 + MAX98357).
    \item Provide ESP32-CAM firmware for MJPEG streaming and pan (stepper) control; document tilt servo limitations imposed by the physical design.
    \item Control a Yeelight bulb (power, brightness, color names, warm/cool white, RGB cycle) from parsed user intent.
    \item Document system architecture, wiring, protocols, and limitations for reproducibility.
\end{itemize}

\section{Thesis outline}

Chapter~2 reviews related classes of systems (voice assistants, smart lighting, ESP32-based vision, social/educational robots, DIY kits) and core technologies (Flask, WebSockets/Socket.IO, Whisper-family STT, LLM tool use, Yeelight LAN control). Chapter~3 presents design and methodology: system overview; physical appearance and CAD; rationale for using \textbf{two} ESP32 modules; electrical interconnect and GPIO tables for both boards; web interface design grounded in the actual CSS/HTML/JS implementation; back-end services and API surface; voice pipeline; firmware behaviour; security notes. Chapter~4 discusses outcomes, qualitative evaluation, the reference Yeelight bulb, and limitations. Chapter~5 concludes and suggests future work. Appendix~A collects protocol summaries and pin quick-reference material.

\section{Acronyms and abbreviations}

\begin{table}[htbp]
\centering
\caption{Acronyms used in this thesis.}
\label{tab:acronyms}
\begin{tabular}{@{}ll@{}}
\toprule
\textbf{Acronym} & \textbf{Meaning} \\
\midrule
API & Application Programming Interface \\
ARIA & Name of the assistant product in this project \\
CAGR & Compound Annual Growth Rate \\
HTTP & Hypertext Transfer Protocol \\
I2S & Inter-IC Sound (digital audio bus) \\
JSON & JavaScript Object Notation \\
LAN & Local Area Network \\
LLM & Large Language Model \\
MJPEG & Motion JPEG (multipart JPEG stream) \\
OAuth & Open Authorization (protocol) \\
PCM & Pulse-Code Modulation \\
RAG & Retrieval-Augmented Generation \\
REST & Representational State Transfer \\
SPI & Serial Peripheral Interface \\
SQL & Structured Query Language \\
STT & Speech-to-Text \\
TLS & Transport Layer Security \\
TTS & Text-to-Speech \\
UDP & User Datagram Protocol \\
UI & User Interface \\
URL & Uniform Resource Locator \\
Wi-Fi & Wireless local networking (IEEE 802.11 family) \\
\bottomrule
\end{tabular}
\end{table}

% -----------------------------------------------
\chapter{Literature review and related work}

\section{Voice assistants and smart-home platforms}

General-purpose voice assistants (e.g., Amazon Alexa, Google Assistant, Apple Siri) couple wake-word detection, cloud speech recognition, and template or LLM-backed dialog managers with vendor-specific smart-home frameworks \cite{flask}. They achieve low perceived latency through optimized end-to-end pipelines and proprietary hardware, but customization is bounded by SDK rules and account linking. ARIA inverts part of this trade-off: the user runs the orchestration server locally and treats ESP32 boards as peripherals, which is closer to the \textit{edge + cloud hybrid} pattern common in IoT research, where heavy models may run on a PC or server while microcontrollers handle sensing and actuation \cite{esp-idf}.

\section{Speech recognition and synthesis}

Modern STT builds on large-scale supervised training; Whisper and its re-implementations demonstrated that a single model can generalize across languages and noise conditions \cite{whisper}. The \textit{faster-whisper} project reuses Whisper weights with efficient runtimes suitable for CPU deployment \cite{faster-whisper}. For TTS without hosting proprietary voice servers, Microsoft's Edge browser voices are exposed programmatically via the \textit{edge-tts} package \cite{edge-tts}, which fits a student project needing multiple locales. Optional wake-word detection is implemented with \textit{openWakeWord} \cite{openwakeword}.

\section{Retrieval and tool-augmented dialogue}

Large language models can be grounded with retrieved documents or structured tool outputs; retrieval-augmented generation (RAG) is a standard pattern for knowledge-intensive tasks \cite{rag}. In ARIA, ``retrieval'' appears pragmatically as SQLite-backed email caches and API results (weather, YouTube search) injected into prompts rather than a full vector database, but the design could be extended.

\section{Positioning}

Table~\ref{tab:comparison_detailed} contrasts ARIA with representative categories. It is not exhaustive but clarifies the niche: \textbf{open-source} stack, \textbf{PC-hosted} intelligence, \textbf{digital I2S audio}, and a \textbf{physical camera head} with pan actuation.

\begin{table}[htbp]
\centering
\renewcommand{\arraystretch}{1.35}
\small
\begin{tabularx}{\textwidth}{|>{\raggedright\arraybackslash}m{3.1cm}|>{\raggedright\arraybackslash}X|>{\raggedright\arraybackslash}X|>{\raggedright\arraybackslash}X|>{\raggedright\arraybackslash}X|}
\hline
\rowcolor{tableblue}
\color{white}\textbf{Criterion} &
\color{white}\textbf{ARIA} &
\color{white}\textbf{Smart speakers} &
\color{white}\textbf{Social robots} &
\color{white}\textbf{DIY kits} \\ \hline

\textbf{LLM depth} & High (Gemini API) & Medium (vendor models / skills) & Low--medium & Often none \\ \hline
\textbf{Physical sensing} & ESP32-CAM + pan stepper & None & Rich (mobility, expression) & Varies \\ \hline
\textbf{Openness} & Open repo, user-hosted & Closed ecosystem & Proprietary & Open hardware \\ \hline
\textbf{Audio path} & I2S INMP441 / MAX98357 & Integrated mics/speakers & Integrated & Often analog buzzers \\ \hline
\end{tabularx}
\caption{High-level comparison of ARIA with other product classes.}
\label{tab:comparison_detailed}
\end{table}

\section{Smart lighting and local control}

Wi-Fi RGBWW bulbs are a common entry point into smart homes because installation is trivial (E27/E26 screw base) and brightness or color can signal assistant state. Yeelight devices historically exposed a \textbf{local LAN control} TCP protocol in addition to cloud app pairing, which is attractive for low-latency automation without round-trips to a vendor cloud \cite{yeelight}. The Python \texttt{yeelight} package wraps this protocol and enables scripts (or Flask handlers) to set power, brightness, color temperature, and RGB. ARIA uses that approach so lamp commands remain on the LAN when the bulb is reachable, while chat and voice layers can still describe lighting changes in natural language.

\section{Embedded vision on ESP32}

The ESP32-CAM module (typically AI-Thinker layout) pairs an ESP32 SoC with an OV2640 (or similar) sensor using a \textbf{parallel camera interface}. Numerous tutorials and open-source sketches demonstrate MJPEG streaming over HTTP \cite{esp-idf}. ARIA extends the minimal web UI with stepper-driven pan and servo tilt hooks. From a systems perspective, continuous frame capture and Wi-Fi transmission already consume GPIO bandwidth, DMA, and CPU time; adding unrelated high-rate peripherals on the same MCU increases scheduling pressure. This motivates the architectural decision analyzed in Section~\ref{sec:dual-esp}.

\section{Open orchestration platforms (context)}

Projects such as Home Assistant aggregate hundreds of vendor integrations through a single hub. ARIA is \textit{not} a drop-in replacement for such platforms; instead it is a \textbf{focused diploma implementation} that prioritizes conversational AI, a custom physical device, and a student-maintainable codebase. The comparison is still useful: it clarifies that ARIA's value lies in transparency of the pipeline (Flask routes, prompt text, firmware pins) rather than breadth of supported SKUs.

% -----------------------------------------------
\chapter{Design and methodology}

\section{System overview}

Figure~\ref{fig:system-arch} summarizes the logical architecture. A desktop or laptop on the LAN runs the Flask process and optional GPU/CPU models. Browser clients talk HTTP and Socket.IO. The ESP32 audio module streams microphone PCM over UDP to the PC and plays back TTS PCM received on a second UDP port. The ESP32-CAM serves MJPEG video and accepts HTTP commands for pan and (intended) tilt. A Yeelight bulb is controlled over the LAN using the manufacturer's local LAN protocol (Python \texttt{yeelight} library \cite{yeelight}).

\begin{figure}[htbp]
\centering
\begin{tikzpicture}[
  font=\small,
  box/.style={draw, rounded corners=2pt, align=center, minimum width=2.6cm, minimum height=0.85cm},
  arrow/.style={-{Latex[length=2.2mm]}, thick}
]
\node[box, fill=tableblue!25] (browser) {Web browser};
\node[box, right=1.1cm of browser, fill=tableblue!40, minimum width=3.2cm] (flask) {Flask +\\Socket.IO};
\node[box, right=1.1cm of flask, fill=gray!15, minimum width=3.4cm] (cloud) {External APIs\\(Gemini, Gmail,\\YouTube, weather)};
\draw[arrow] (browser) -- node[above, font=\scriptsize] {HTTP / WS} (flask);
\draw[arrow] (flask) -- node[above, font=\scriptsize] {HTTPS} (cloud);

\node[box, below=1.15cm of browser, fill=green!12] (cam) {ESP32-CAM\\{\scriptsize HTTP :80 / :81}};
\node[box, below=1.15cm of flask, fill=green!12] (audio) {ESP32-WROOM\\{\scriptsize UDP mic/spk}};
\node[box, below=1.15cm of cloud, fill=yellow!20] (bulb) {Yeelight\\{\scriptsize TCP LAN}};

\draw[arrow] (flask.south west) -- ++(0,-0.35) -| node[pos=0.22, above, font=\scriptsize] {LAN} (cam.north);
\draw[arrow] (flask) -- node[right, font=\scriptsize, xshift=2pt] {LAN} (audio);
\draw[arrow] (flask.south east) -- ++(0,-0.35) -| (bulb.north);
\end{tikzpicture}
\caption{Logical deployment of ARIA on a home LAN (browser, gateway, peripherals, cloud APIs).}
\label{fig:system-arch}
\end{figure}

\textbf{Repository and reproducibility.} Source code and firmware sketches are maintained in a Git repository; the public reference cited in development is \url{https://github.com/TheSuxarick/aria}.

\section{Physical and mechanical design}

\subsection{Design language and appearance}

The physical prototype targets a \textbf{compact, monolithic silhouette}: two stacked spherical segments (a larger lower ``body'' and a smaller upper ``head''), inspired by popular culture robot aesthetics but implemented as an original CAD model. The exterior is printed in \textbf{matte black} filament; layer lines and slight faceting from the mesh tessellation remain visible and read as an honest manufacturing trace rather than a flaw---similar to consumer FDM parts where texture signals hand-built iteration.

The head segment carries the \textbf{camera aperture}: a circular front housing with the ESP32-CAM lens and auxiliary LED activity. In the assembled unit the indicator glow reads as a \textbf{``red eye''} in low light, which doubles as an intuitive ``device live'' cue for demos and photography. Figure~\ref{fig:final-look} shows the finished device next to the high-current 5~V brick used to feed regulators, amplifiers, and motors.

\begin{figure}[p]
\centering
\includegraphics[width=0.72\textwidth]{final_physical_look.png}
\caption{Assembled ARIA prototype (matte black 3D-printed enclosure, camera ``eye'' with status lighting) and external 5~V DC supply.}
\label{fig:final-look}
\end{figure}

\subsection{CAD geometry and internal layout}

Before printing, the shell was modeled in CAD software. Figure~\ref{fig:3d-front} and Figure~\ref{fig:3d-back} show orthographic-style views of the outer surfaces: the head dome, the front camera recess, and the lower sphere volume that hides audio electronics and cabling exits. Figure~\ref{fig:3d-xray} presents a \textbf{see-through (wireframe)} view that makes internal intent explicit: a vertical neck region reserved for pan mechanics, horizontal shelves inside the body for the ESP32-WROOM stack, INMP441, MAX98357, and power distribution, and the rectangular cutout in the head aligned with the camera module.

\begin{figure}[p]
\centering
\includegraphics[width=0.62\textwidth]{3dmodel_front.png}
\caption{CAD model: front aspect showing head dome and camera recess.}
\label{fig:3d-front}
\end{figure}

\begin{figure}[p]
\centering
\includegraphics[width=0.62\textwidth]{3dmodel_bottom_back.png}
\caption{CAD model: bottom and rear regions (cable exit and base footprint).}
\label{fig:3d-back}
\end{figure}

\begin{figure}[p]
\centering
\includegraphics[width=0.78\textwidth]{3dmodel_see_through.png}
\caption{Transparent/wireframe CAD view illustrating internal volumes: camera pocket in the head, neck/pan region, and body shelves for audio MCU and power.}
\label{fig:3d-xray}
\end{figure}

\subsection{Power and electromechanical notes}

The prototype is powered from a \textbf{5~V DC supply rated for about 10~A}. This margin covers inrush and simultaneous load from two ESP32 modules, the class-D amplifier peaks, stepper coils, and conversion losses. Each downstream rail (3.3~V for logic and microphone, 5~V for the speaker amplifier where applicable) should be fused or protected as appropriate for the final wiring harness.

\textbf{Mechanical limitation.} Firmware includes a servo-driven tilt axis, but the printed assembly and tolerances did not yield reliable tilt travel; operationally, \textbf{horizontal pan} (stepper) is the dependable repositioning degree of freedom for the camera.

\section{Rationale for two ESP32 modules}
\label{sec:dual-esp}

The product uses \textbf{one ESP32-CAM} and \textbf{one ESP32-WROOM-class dev board} instead of a single MCU. The honest project motivation included \textbf{GPIO exhaustion}: the AI-Thinker camera module dedicates the majority of useful pins to the parallel camera data bus, I$^2$C SCCB, synchronization lines, and PSRAM-friendly routing. Adding \textbf{two independent I2S peripherals} (microphone RX and speaker TX) each requires BCLK, LRCLK/WSEL, and data lines, plus stable DMA service---overlapping those resources with continuous camera DMA and dual HTTP servers increases interrupt latency and debugging difficulty.

Stated more formally, the split follows good embedded practice:

\begin{itemize}
    \item \textbf{Pin budget and contention.} Tables~\ref{tab:cam-gpio-camera}--\ref{tab:cam-gpio-act} show that the OV2640 parallel interface alone occupies many GPIOs on the CAM board; remaining free pins are insufficient for a clean dual-I2S audio layout alongside stepper and servo outputs without resorting to risky pin sharing or stripping camera features.
    \item \textbf{Real-time isolation.} MJPEG encoding and Wi-Fi TX on the camera MCU run continuously during streaming. Full-duplex audio on the same chip would share the RF stack and cache behavior; separating audio to a second ESP32 isolates dropouts: a camera watchdog reboot does not necessarily mute the microphone bridge.
    \item \textbf{Firmware simplicity.} Each sketch stays small and testable: \texttt{esp-32-cam.ino} handles vision and motion; \texttt{esp-32-wroom.ino} handles only I2S and UDP. Students can flash and bench-test audio without touching camera registers, and vice versa.
    \item \textbf{Power and mechanical packaging.} Physically, the head hosts the camera board while the body hosts audio---matching the CAD intent in Figure~\ref{fig:3d-xray}.
\end{itemize}

A single-chip solution might be attempted on a custom PCB with a larger ESP32-S3 module, external PSRAM, careful I2S pinmux, and a unified FreeRTOS design; that was outside diploma time and fabrication budget.

\section{Electrical interconnect}

\subsection{ESP32-CAM module}

The firmware selects \texttt{CAMERA\_MODEL\_AI\_THINKER}. Tables~\ref{tab:cam-gpio-camera} and \ref{tab:cam-gpio-act} list GPIO assignments as defined in \texttt{esp-32-cam.ino}. Many lines are dictated by the vendor reference design; the stepper and servo lines are project-specific actuator wiring. \textbf{Note:} GPIO2 is a strapping pin on ESP32; using it for a stepper phase requires care at reset---the firmware assumes a standard hobby-driver setup where this has been validated.

\begin{table}[htbp]
\centering
\caption{ESP32-CAM (AI-Thinker): parallel camera interface and control signals (from firmware).}
\label{tab:cam-gpio-camera}
\begin{tabular}{@{}lll@{}}
\toprule
\textbf{Function} & \textbf{Signal / role} & \textbf{GPIO} \\
\midrule
Camera power down & PWDN & 32 \\
Camera reset & RESET & not used ($-1$) \\
External clock & XCLK & 0 \\
SCCB data / clock & SIOD / SIOC & 26 / 27 \\
Pixel data & Y2 \ldots Y9 & 5, 18, 19, 21, 36, 39, 34, 35 \\
Frame sync & VSYNC & 25 \\
Line valid & HREF & 23 \\
Pixel clock & PCLK & 22 \\
\bottomrule
\end{tabular}
\end{table}

\begin{table}[htbp]
\centering
\caption{ESP32-CAM (AI-Thinker): actuators and motion (project wiring in firmware).}
\label{tab:cam-gpio-act}
\begin{tabular}{@{}lll@{}}
\toprule
\textbf{Function} & \textbf{Description} & \textbf{GPIO} \\
\midrule
Stepper winding A1 & \texttt{STEPPER\_PIN1} & 13 \\
Stepper winding A2 & \texttt{STEPPER\_PIN2} & 15 \\
Stepper winding B1 & \texttt{STEPPER\_PIN3} & 14 \\
Stepper winding B2 & \texttt{STEPPER\_PIN4} & 2 \\
Tilt servo PWM & \texttt{SERVO\_1} & 12 \\
\bottomrule
\end{tabular}
\end{table}

Figure~\ref{fig:cam-wiring} provides the project wiring drawing for the camera assembly (to be aligned with the tables above).

\begin{figure}[htbp]
\centering
\includegraphics[width=0.92\textwidth]{esp32-cam_connection_sheme.png}
\caption{Wiring diagram for the ESP32-CAM assembly (camera module, pan stepper, tilt servo signal).}
\label{fig:cam-wiring}
\end{figure}

\subsection{ESP32-WROOM audio module}

The audio bridge uses two I2S peripherals: RX for the INMP441 and TX for the MAX98357. Table~\ref{tab:audio-pins} lists GPIOs as implemented in repository firmware (\texttt{esp-32-wroom.ino}). GPIO~34 is input-only and is used for the microphone data line per the prototype wiring.

\begin{table}[htbp]
\centering
\caption{ESP32 (WROOM) I2S pin assignment in repository firmware.}
\label{tab:audio-pins}
\begin{tabular}{lll}
\toprule
\textbf{Function} & \textbf{Signal} & \textbf{GPIO} \\
\midrule
INMP441 & Bit clock (SCK / BCLK) & 26 \\
INMP441 & Word select (WS / LRCLK) & 25 \\
INMP441 & Serial data (SD) & 34 \\
MAX98357 & Bit clock (BCLK) & 12 \\
MAX98357 & Left-right clock (LRC) & 14 \\
MAX98357 & Data in (DIN) & 27 \\
\bottomrule
\end{tabular}
\end{table}

\begin{figure}[htbp]
\centering
\includegraphics[width=0.92\textwidth]{esp-32-wroom_connection_scheme.png}
\caption{Wiring diagram for the ESP32-WROOM audio node (INMP441 + MAX98357).}
\label{fig:wroom-wiring}
\end{figure}

\textbf{Audio protocol.} Both directions use \textbf{16-bit signed PCM, mono, 16~kHz}, little-endian, over UDP. The microphone path sends datagrams to \texttt{PC\_IP:12345}; the PC sends speaker frames to the ESP on port~\texttt{12346}.

\section{Software methodology}

Development followed \textbf{iterative prototyping with hardware--software co-design}: camera and audio firmware were tested against early Flask endpoints before the dashboard was finalized; Yeelight and Gmail features were added as independent modules behind HTTP routes to reduce coupling.

\section{Web application design}

\subsection{Information architecture}

The interface is implemented as a \textbf{single-page application shell}: the server renders \texttt{index.html} once through Jinja, and navigation swaps content panels in the DOM without full page reloads. A fixed \textbf{sidebar} (width 240\,px in CSS variables) lists four primary destinations---\textit{Dashboard}, \textit{Camera}, \textit{Functions}, \textit{Settings}---each bound to a \texttt{data-page} attribute consumed by \texttt{app.js}. A \textbf{top header} mirrors the current title, exposes a theme toggle, and shows a compact \textit{Server: Online} badge that tracks reachability of the Flask host.

This structure mirrors common enterprise dashboards: persistent global navigation, contextual main region, and status indicators that reassure the user that automation commands will reach the back end. It also maps cleanly onto the hardware split: lamp and chat on the dashboard; MJPEG and pan on the camera page; long-horizon information (forecast, inbox) on functions; audio routing and LLM persona on settings.

\subsection{Visual design language and theming}

The stylesheet (\texttt{style.css}) defines a \textbf{token-based palette} using CSS custom properties on \texttt{:root}. The default (dark) theme anchors the background in slate blues (\texttt{\#0f172a}, \texttt{\#1e293b}) rather than pure black, reducing halation on LCD panels. The accent is a saturated \textbf{emerald green} (\texttt{\#10b981}) used for active nav pills, primary buttons, scrollbar thumbs, and success states---a deliberate choice to evoke ``go / active / home IoT'' without relying on generic corporate blue.

A companion \textbf{light theme} is activated by setting \texttt{data-theme="light"} on the document root; the same variable names remap to white cards, softer shadows, and slightly darker accent green (\texttt{\#059669}) for WCAG-friendly contrast on white. Toggle persistence (localStorage) lets operators match room lighting during demos.

Typography pairs \textbf{Inter} (loaded from Google Fonts in the template) with \textbf{Material Icons Round} for wayfinding. Cards use 12--16\,px corner radii, layered shadows, and subtle radial background gradients behind the main canvas so empty states still feel intentional.

\subsection{Dashboard view}

The landing view aggregates \textbf{three interaction modes} aligned with project goals: (1)~\textit{Lamp control}---Yeelight power, brightness slider (default 80\%), warm/day/cool white presets, eight named color swatches, and an RGB cycle mode; (2)~\textit{Voice robot}---push-to-talk style control that records through the selected audio source (ESP32 or browser) and plays replies; (3)~\textit{AI Interaction Chat}---textual Gemini dialogue with scrollable history and a clear-history affordance.

Figure~\ref{fig:ui-dash} shows the dashboard in light theme: the lamp card exposes state (\textsc{off}/\textsc{on}), the robot card prompts the user through the call workflow, and the chat card illustrates multilingual use (e.g., Russian weather queries) while the chrome stays English for consistency with the development team's working language.

\begin{figure}[p]
\centering
\includegraphics[width=0.92\textwidth]{ui_dashboard.png}
\caption{Dashboard view: Yeelight control, voice robot card, and AI chat with multilingual messages.}
\label{fig:ui-dash}
\end{figure}

\subsection{Camera and telepresence controls}

The camera page embeds the ESP32-CAM MJPEG stream in a large preview region. When the stream URL is wrong or the module is offline, the UI shows a \textbf{graceful empty state} (camera glyph, explanatory text, \textit{Retry}) and raises a toast (\textit{Camera not found}) so users distinguish software errors from network errors.

Below the preview, a \textbf{D-pad} issues HTTP GET requests to \texttt{/action?go=...} on the camera module (Up/Down for servo tilt, Left/Right for stepper pan). Adjacent toggles control microphone listen path and mute state, linking the visual subsystem with the audio bridge managed on the PC. Figure~\ref{fig:ui-cam} captures this layout; the illustrated error state documents real-world debugging sessions where DHCP leases changed.

\begin{figure}[p]
\centering
\includegraphics[width=0.92\textwidth]{ui_camera.png}
\caption{Camera view: stream placeholder with retry, pan/tilt pad, and audio controls; toast on failure.}
\label{fig:ui-cam}
\end{figure}

\subsection{Functions: weather and email}

The \textit{Functions} page foregrounds \textbf{information retrieval}. A horizontally scrollable \textbf{forecast strip} renders per-interval cards (icon, temperature, wind or humidity) fed by the same OpenWeather-backed endpoints used in chat context. Below it, the \textbf{Email} card lists cached Gmail messages (avatar initials, sender, subject, snippet, relative timestamps) with a \textit{Sync} action. Mixed Cyrillic timestamps illustrate locale-aware formatting in the wild.

Figure~\ref{fig:ui-func} provides a representative screenshot.

\begin{figure}[p]
\centering
\includegraphics[width=0.92\textwidth]{ui_functions.png}
\caption{Functions view: scrollable weather timeline and Gmail list with sync control.}
\label{fig:ui-func}
\end{figure}

\subsection{Settings and personalization}

Settings concentrate \textbf{non-secret preferences}. \textit{AI Personality} offers eight presets (Default, Chill, Bro, Angry, Formal, Pirate, Sassy, Nerd) mapped to alternative system prompts in \texttt{app.py}. \textit{Language} selects UI copy among EN/RU/KZ, reflecting the Kazakhstan deployment context. \textit{Audio source} toggles between \textbf{ESP32} I2S bridge and \textbf{This Device} (browser mic)---critical for development on laptops without hardware. \textit{Wake word} enables ``Computer'' hotword detection via \textit{openWakeWord}. An API key summary card communicates how many Gemini keys were loaded from environment variables without printing secrets. Figure~\ref{fig:ui-set} shows the layout.

\begin{figure}[p]
\centering
\includegraphics[width=0.92\textwidth]{ui_settings.png}
\caption{Settings view: personality presets, language, audio source, wake word, API key status.}
\label{fig:ui-set}
\end{figure}

\subsection{Client technologies and patterns}

\begin{itemize}
    \item \textbf{HTML5/CSS3} --- Flex layouts, CSS variables, responsive sidebar collapse, reduced-motion-friendly transitions (\texttt{0.22s cubic-bezier}).
    \item \textbf{Vanilla JavaScript} --- Modular functions in \texttt{app.js} for lamp state polling, chat transcript rendering with HTML escaping, music intent parsing (multi-language trigger phrases), and Socket.IO namespaces.
    \item \textbf{Progressive enhancement} --- If WebSocket audio is blocked, text chat and REST features remain usable.
    \item \textbf{Security-minded DOM hygiene} --- User and model messages pass through \texttt{escapeHtml} before insertion to mitigate XSS in chat rendering.
\end{itemize}

\section{Back-end architecture and HTTP API}

\subsection{Server stack}

The server is implemented in \textbf{Python~3} using \textbf{Flask} \cite{flask} for HTTP and \textbf{Flask-SocketIO} \cite{flask-socketio} (Socket.IO semantics \cite{socketio}) for streaming events such as microphone audio to the browser and robot status. SQL persistence uses \textbf{Flask-SQLAlchemy} \cite{flask-sqlalchemy} over \textbf{SQLite} \cite{sqlite}. Gmail access follows OAuth~2.0 \cite{rfc6749, google-oauth, gmail-api}. The LLM endpoint uses the \textbf{Gemini} HTTP API \cite{gemini-api}. Weather data uses OpenWeather-compatible requests \cite{openweather}; music search uses the YouTube Data API \cite{youtube-api}. Rate limiting is available via Flask-Limiter \cite{flask-limiter}. Secrets load from \texttt{.env} via \texttt{python-dotenv} \cite{python-dotenv}.

\subsection{Representative routes}

Table~\ref{tab:api-routes} documents the main JSON/REST surfaces used by the dashboard. It is representative rather than exhaustive; OAuth callback routes and static file serving are omitted for brevity.

\begin{table}[htbp]
\centering
\caption{Selected Flask routes in ARIA (representative).}
\label{tab:api-routes}
\small
\begin{tabularx}{\textwidth}{@{}p{0.26\textwidth}p{0.1\textwidth}X@{}}
\toprule
\textbf{Path} & \textbf{Verb} & \textbf{Role} \\
\midrule
\texttt{/api/chat} & POST & Send user text; returns model reply; persists rows in \texttt{chat\_messages}. \\
\texttt{/api/chat/history} & GET/DELETE & Read or clear stored messages for a session id. \\
\texttt{/api/weather} & GET & Current weather JSON for default or requested city. \\
\texttt{/api/forecast} & GET & Multi-step forecast for UI cards. \\
\texttt{/api/play-music} & POST & YouTube search given query; returns video id for embed player. \\
\texttt{/api/lamp-state} & GET & Poll synthesized lamp state for dashboard badges. \\
\texttt{/api/quick-action} & POST & Shortcut actions from UI buttons. \\
\texttt{/api/settings} & GET/POST & Read/update model name, UI language, personality, wake word flags. \\
\texttt{/api/memory} & GET/POST/DELETE & Lightweight long-lived text snippets for prompts. \\
\texttt{/api/server-status} & GET & Health ping used by badges. \\
\bottomrule
\end{tabularx}
\end{table}

Socket.IO namespace \texttt{/audio} carries binary PCM chunks and textual robot state events between browser and server when voice features are enabled.

\subsection{Natural-language lamp control}

Beyond the dashboard sliders, \texttt{app.py} implements \textbf{keyword routers} for lamp commands inside chat utterances. English, Russian, and Kazakh stems map to power toggles, brightness changes, warm/cool white, named RGB presets, and a rainbow cycle. This dual surface (GUI + language) is important for demos where a visitor says ``turn the light blue'' while another operator uses the palette.

\section{Data model}

Key tables (see \texttt{models.py}): \texttt{ChatMessage} for persisted dialogue lines keyed by session; \texttt{User} and \texttt{Session} for optional local accounts; \texttt{GmailAccount} for OAuth tokens; \texttt{EmailMessage} for cached inbox snippets shown to the model.

\begin{lstlisting}[language=Python, caption={Excerpt: chat message model (conceptual).}, label={lst:model}]
class ChatMessage(db.Model):
    __tablename__ = "chat_messages"
    id = db.Column(db.Integer, primary_key=True)
    session_id = db.Column(db.String(64), nullable=False, index=True)
    role = db.Column(db.String(16), nullable=False)
    text = db.Column(db.Text, nullable=False)
\end{lstlisting}

\section{Voice and robot pipeline}

Figure~\ref{fig:voice-flow} sketches the audio data path. The ESP32 captures 16-bit PCM at 16~kHz, packs samples into UDP datagrams, and sends them to the PC listener on port~\texttt{12345}. A background thread in \texttt{app.py} demultiplexes packets, optionally runs \textit{openWakeWord} \cite{openwakeword} frames, buffers speech until silence or timeout, and calls \textit{faster-whisper} \cite{faster-whisper, whisper} (typically \texttt{small} on CPU \texttt{int8}, with preload at startup). The recognized string is injected into the same Gemini orchestration path as typed chat. Responses are rendered to speech via \textit{edge-tts} \cite{edge-tts}; PCM is sent back to the ESP32 speaker task on port~\texttt{12346}. Browser mode skips UDP and uses WebSocket audio where supported. \texttt{pydub} assists when wave containers are needed \cite{pydub}.

\begin{figure}[htbp]
\centering
\begin{tikzpicture}[font=\small, node distance=0.85cm,
  blk/.style={draw, rounded corners, align=center, minimum height=0.75cm, minimum width=2.2cm}]
\node[blk, fill=green!15] (mic) {INMP441};
\node[blk, right=of mic, fill=green!15] (esp) {ESP32\\UDP out};
\node[blk, right=of esp, fill=blue!12] (py) {Flask bridge\\+ Whisper};
\node[blk, right=of py, fill=purple!15] (llm) {Gemini};
\node[blk, below=of py, fill=orange!15] (tts) {edge-tts};
\node[blk, left=of tts, fill=green!15] (spk) {MAX98357};
\draw[-{Latex}] (mic) -- (esp);
\draw[-{Latex}] (esp) -- node[above, font=\scriptsize] {12345} (py);
\draw[-{Latex}] (py) -- (llm);
\draw[-{Latex}] (llm) |- (tts);
\draw[-{Latex}] (tts) -- node[below, font=\scriptsize] {12346} (spk);
\end{tikzpicture}
\caption{Simplified voice pipeline: microphone $\rightarrow$ ESP32 UDP $\rightarrow$ STT/LLM $\rightarrow$ TTS $\rightarrow$ speaker.}
\label{fig:voice-flow}
\end{figure}

\section{Device firmware}

\textbf{ESP32-CAM} disables brown-out noise during boot, configures the OV2640 over SCCB, allocates frame buffers with PSRAM when available, starts two HTTP listeners (see Appendix~\ref{app:a}), spawns a FreeRTOS queue consumer for stepper commands, and maps UI directions to signed step counts. Tilt uses a 50\,Hz servo signal on GPIO~12.

\textbf{ESP32-WROOM} installs separate I2S drivers for RX and TX, maps INMP441 32-bit frames to int16 PCM, applies configurable speaker gain before \texttt{i2s\_write}, and pins microphone and speaker loops to different cores to reduce dropouts.

\textbf{Application launcher.} The repository ships \texttt{run\_server.py} at the project root; it inserts \texttt{ARIA website} on \texttt{sys.path}, changes the working directory, and invokes \texttt{socketio.run} on \texttt{0.0.0.0:5000} for LAN access from phones and ESP32 devices.

\section{Security and privacy notes}

The prototype targets a \textbf{trusted home or lab LAN}. Assets at risk include: (i)~OAuth refresh tokens for Gmail stored in SQLite; (ii)~Gemini and YouTube API keys in environment variables; (iii)~live microphone audio traversing UDP and WebSockets; (iv)~MJPEG video URLs that may be guessable if IP addresses leak.

Mitigations implemented in code include: loading secrets from \texttt{.env} rather than hard-coding \cite{python-dotenv}; HTTP-only session cookies with SameSite lax defaults for OAuth compatibility; redirecting \texttt{127.0.0.1} to \texttt{localhost} to stabilize OAuth redirect URI matching; and HTML escaping of chat content before DOM insertion.

Residual risks and hardening steps for a production deployment would include: HTTPS reverse proxy (nginx/Caddy) with valid certificates; pinning CORS and Socket.IO origins to known UI hosts; encrypting the SQLite file at rest if tokens are sensitive; VLAN or guest-network isolation for IoT devices; and periodic key rotation with monitoring of API quotas.

% -----------------------------------------------
\chapter{Results and discussion}

\section{Functional outcome}

The integrated system supports: browser-based chat with Gemini; reading and summarizing cached Gmail messages when OAuth is configured; weather and forecast queries; starting YouTube playback from natural language; toggling and recoloring a Yeelight bulb via keyword rules in multiple languages; viewing the ESP32-CAM stream inside the dashboard; and, when the PC pipeline is active, two-way voice through the I2S peripherals.

\section{Demonstration scenarios}

To validate end-to-end behaviour during development, the team repeated several \textbf{scenario scripts} (informal acceptance tests):

\begin{enumerate}
    \item \textbf{Morning briefing.} User asks for weather in Russian in chat; ARIA returns structured conditions; the same intent works from voice when the ESP32 audio bridge is online.
    \item \textbf{Lighting demo.} Operator sets warm white at 80\% from the dashboard, then asks in chat to switch to blue and dim; lamp state badge updates consistently.
    \item \textbf{Inbox triage.} After Gmail OAuth, user requests a summary of the latest retail newsletter; model consumes cached \texttt{EmailMessage} rows without re-fetching the full thread from Google for every token.
    \item \textbf{Telepresence.} User opens the camera page, verifies MJPEG latency on LAN, pans left/right to track a subject; tilt is attempted but may stall depending on mechanical clearance.
    \item \textbf{Fallbacks.} Disconnect ESP32-CAM power; UI shows empty state and toast (Figure~\ref{fig:ui-cam}); disconnect Internet; chat degrades with explicit error rather than silent failure where implemented.
\end{enumerate}

These scenarios do not constitute a formal user study with $N$ participants; they are engineering walkthroughs suitable for reproducibility in defense sessions.

\section{Qualitative observations}

\textbf{Latency.} Text chat latency is dominated by Gemini round-trip time. Voice adds STT and TTS queues; perceived gap is acceptable on a mid-range laptop CPU with \texttt{small} Whisper but grows if antivirus software scans model files on first load.

\textbf{Stability.} Flask-SocketIO over threading mode proved sufficient for a single-household demo; concurrent multi-user load was not a goal. UDP audio occasionally clipped when Wi-Fi RSSI dropped below roughly $-75$\,dBm in informal tests---quantitative PER was not measured.

\textbf{UX.} The mint accent and card layout received positive informal feedback for readability; the camera error state was considered clearer than a broken image icon alone.

\section{BOM and cost-oriented summary}

Table~\ref{tab:bom} lists indicative major hardware (retail prices vary by region and vendor). The table supports claims that the physical prototype stays within a student budget while still demonstrating multimodal integration.

\begin{table}[htbp]
\centering
\caption{Indicative bill of materials (major components only).}
\label{tab:bom}
\small
\begin{tabular}{@{}p{0.42\textwidth}p{0.48\textwidth}@{}}
\toprule
\textbf{Component} & \textbf{Notes} \\
\midrule
ESP32-CAM (AI-Thinker) & Camera + Wi-Fi SoC on one board. \\
ESP32-WROOM-32 DevKit class board & Hosts dual I2S (mic + DAC). \\
INMP441 breakout & Digital MEMS microphone, I2S. \\
MAX98357A + speaker & I2S class-D amplifier module. \\
Stepper + driver + mechanical prints & Pan axis; filament cost minor. \\
Yeelight YLDP005 & Validated smart bulb (Table~\ref{tab:yeelight-bulb}). \\
5~V / 10~A DC supply & Headroom for Wi-Fi peaks and audio. \\
PC (existing) & Hosts Flask, Whisper, TTS, browsers. \\
\bottomrule
\end{tabular}
\end{table}

\subsection{Yeelight lamp used in tests}

Lighting control was validated on a consumer Yeelight Wi-Fi bulb with local LAN API support. Table~\ref{tab:yeelight-bulb} lists the unit exercised during integration testing; other Yeelight models that expose the same LAN protocol should behave similarly, but only this SKU was confirmed in the lab.

\begin{table}[htbp]
\centering
\caption{Yeelight smart bulb used for ARIA lighting tests.}
\label{tab:yeelight-bulb}
\begin{tabular}{@{}ll@{}}
\toprule
\textbf{Attribute} & \textbf{Value} \\
\midrule
Product name & Smart LED Bulb W3 (Multicolor) \\
Manufacturer model & YLDP005 \\
Mains supply & 100--240~V, 50/60~Hz; rated current 0.07~A \\
Optical output & 900~lm nominal \\
Color temperature & 1700--6500~K (white range); multicolor (RGB) modes \\
Rated lamp power & 8~W \\
\bottomrule
\end{tabular}
\end{table}

\section{Limitations}

\begin{itemize}
    \item \textbf{Mechanics:} tilt servo travel was limited by the printed assembly; pan is the dependable axis.
    \item \textbf{Network:} devices assume stable LAN connectivity; UDP audio can exhibit dropouts under congestion unless packet sizes and Wi-Fi placement are tuned.
    \item \textbf{Cloud dependence:} Gemini, Gmail, YouTube, and weather APIs require Internet and suitable API keys; only SQLite state and LAN bulb control remain purely local.
    \item \textbf{Wake word:} accuracy depends on model choice, microphone SNR, and CPU load; it is optional in settings.
    \item \textbf{Evaluation depth:} no formal usability experiment or benchmark suite; results are qualitative engineering validation.
    \item \textbf{Single-bulb assumption:} lamp discovery uses MAC/ARP heuristics tuned to one fixture; multi-room lighting would need provisioning UX.
\end{itemize}

% -----------------------------------------------
\chapter{Conclusion}

This thesis presented \textbf{ARIA}, a multimodal smart-home assistant that combines a Flask-hosted web dashboard, Google Gemini--based dialogue, local speech processing (\textit{faster-whisper}, \textit{edge-tts}), OAuth-backed Gmail caching, multimedia search, Yeelight LAN control, and a 3D-printed physical embodiment. Two ESP32 microcontrollers partition the problem along \textbf{pin-count, real-time, and packaging boundaries}: one module concentrates on parallel camera IO, MJPEG streaming, and pan motion; the other dedicates two I2S peripherals to full-duplex digital audio bridged to the PC over UDP. That split is not merely an accident of wiring---it reflects how modern embedded products isolate subsystems to keep firmware testable and failures contained.

The web interface demonstrates that a \textbf{vanilla HTML/CSS/JS} stack remains sufficient for polished dashboards when design tokens, theming, and component cards are disciplined. Screenshots of dashboard, camera, functions, and settings views (Chapter~3) document the final UX for examination boards even when the repository evolves.

Future work includes: mechanical redesign to restore reliable tilt; HTTPS deployment and stricter Socket.IO origin policy; provisioning flows for multiple Yeelight devices; on-device wake-word or VAD on the audio ESP32 to reduce PC idle load; optional local LLM inference for offline mode; and deeper quantitative evaluation (WER for STT, camera latency histograms, user studies).

The project shows that a small team can integrate commodity IoT parts, open Python libraries, and cloud LLMs into a coherent assistant without surrendering all data paths to a single vendor appliance---a relevant direction as households in Kazakhstan and abroad demand Russian, Kazakh, and English interfaces in one household system.

\begin{thebibliography}{99}

\bibitem{technavio}
\textbf{Global Smart Home Market 2025--2029 (press summary).} Technavio.
\newblock \url{https://newsroom.technavio.org/smart-homemarket-v2}
\newblock Accessed: 2026-04-12.

\bibitem{flask}
\textbf{Flask Documentation.} Pallets Projects.
\newblock \url{https://flask.palletsprojects.com/en/stable/}
\newblock Accessed: 2026-04-11.

\bibitem{flask-socketio}
\textbf{Flask-SocketIO Documentation.} Miguel Grinberg.
\newblock \url{https://flask-socketio.readthedocs.io/en/latest/}
\newblock Accessed: 2026-04-11.

\bibitem{socketio}
\textbf{Socket.IO Documentation (v4).}
\newblock \url{https://socket.io/docs/v4/}
\newblock Accessed: 2026-04-11.

\bibitem{sqlalchemy}
\textbf{SQLAlchemy 2.0 Documentation.}
\newblock \url{https://docs.sqlalchemy.org/en/20/}
\newblock Accessed: 2026-04-11.

\bibitem{flask-sqlalchemy}
\textbf{Flask-SQLAlchemy Documentation.} Pallets Projects.
\newblock \url{https://flask-sqlalchemy.palletsprojects.com/en/stable/}
\newblock Accessed: 2026-04-11.

\bibitem{sqlite}
\textbf{SQLite Documentation.}
\newblock \url{https://www.sqlite.org/docs.html}
\newblock Accessed: 2026-04-11.

\bibitem{rfc6749}
Hardt, D.
\newblock \textbf{The OAuth 2.0 Authorization Framework.} RFC~6749, IETF, 2012.
\newblock \url{https://datatracker.ietf.org/doc/html/rfc6749}

\bibitem{google-oauth}
\textbf{Using OAuth 2.0 to Access Google APIs.} Google.
\newblock \url{https://developers.google.com/identity/protocols/oauth2}
\newblock Accessed: 2026-04-11.

\bibitem{gmail-api}
\textbf{Gmail API Overview.} Google.
\newblock \url{https://developers.google.com/gmail/api/guides}
\newblock Accessed: 2026-04-11.

\bibitem{gemini-api}
\textbf{Gemini API Documentation.} Google.
\newblock \url{https://ai.google.dev/gemini-api/docs}
\newblock Accessed: 2026-04-11.

\bibitem{openweather}
\textbf{Weather APIs.} OpenWeather.
\newblock \url{https://openweathermap.org/api}
\newblock Accessed: 2026-04-11.

\bibitem{youtube-api}
\textbf{YouTube Data API v3.} Google.
\newblock \url{https://developers.google.com/youtube/v3}
\newblock Accessed: 2026-04-11.

\bibitem{flask-limiter}
\textbf{Flask-Limiter Documentation.} Ali-Akber Saifee.
\newblock \url{https://flask-limiter.readthedocs.io/en/stable/}
\newblock Accessed: 2026-04-11.

\bibitem{whisper}
Radford, A., Kim, J.~W., Xu, T., Brockman, G., McLeavey, C., Sutskever, I.
\newblock \textbf{Robust Speech Recognition via Large-Scale Weak Supervision.} arXiv:2212.04356, 2022.
\newblock \url{https://arxiv.org/abs/2212.04356}

\bibitem{faster-whisper}
\textbf{faster-whisper.} SYSTRAN.
\newblock GitHub repository.
\newblock \url{https://github.com/SYSTRAN/faster-whisper}
\newblock Accessed: 2026-04-11.

\bibitem{edge-tts}
\textbf{edge-tts.} PyPI package.
\newblock \url{https://pypi.org/project/edge-tts/}
\newblock Accessed: 2026-04-11.

\bibitem{openwakeword}
\textbf{openWakeWord.} David Scripka et al.
\newblock \url{https://github.com/dscripka/openWakeWord}
\newblock Accessed: 2026-04-12.

\bibitem{pydub}
\textbf{pydub.} James Robertson.
\newblock \url{https://github.com/jiaaro/pydub}
\newblock Accessed: 2026-04-11.

\bibitem{esp-idf}
\textbf{Get Started --- ESP32 --- ESP-IDF Programming Guide.} Espressif Systems.
\newblock \url{https://docs.espressif.com/projects/esp-idf/en/latest/esp32/get-started/index.html}
\newblock Accessed: 2026-04-11.

\bibitem{yeelight}
\textbf{yeelight (Python library).}
\newblock \url{https://yeelight.readthedocs.io/}
\newblock Accessed: 2026-04-11.

\bibitem{rag}
Lewis, P., et al.
\newblock \textbf{Retrieval-Augmented Generation for Knowledge-Intensive NLP Tasks.} arXiv:2005.11401, 2020.
\newblock \url{https://arxiv.org/abs/2005.11401}

\bibitem{requests}
\textbf{Requests: HTTP for Humans.}
\newblock \url{https://requests.readthedocs.io/en/latest/}
\newblock Accessed: 2026-04-11.

\bibitem{python-dotenv}
\textbf{python-dotenv.} PyPI package.
\newblock \url{https://pypi.org/project/python-dotenv/}
\newblock Accessed: 2026-04-11.

\end{thebibliography}

\appendix

\chapter{Appendix A: Wiring and protocol summary}
\label{app:a}

\section{ESP32-CAM control endpoints}

\begin{itemize}
    \item \textbf{HTTP port 80:} \texttt{/} (minimal HTML UI), \texttt{/action?go=left|right|up|down} (pan/tilt commands).
    \item \textbf{HTTP port 81:} \texttt{/stream} --- MJPEG multipart stream for embedding in the dashboard.
\end{itemize}

\section{ESP32 audio UDP ports}

\begin{itemize}
    \item \textbf{12345} --- microphone PCM from ESP32 to PC (PC listens).
    \item \textbf{12346} --- speaker PCM from PC to ESP32 (ESP listens).
\end{itemize}

\section{Yeelight control}

Bulb commands are issued on the LAN once the lamp IP is resolved (ARP table / environment fallback). Natural-language rules in the server map phrases in English, Russian, and Kazakh to power and color changes.

\subsection{Reference hardware}

The stack was tested with a \textbf{Yeelight Smart LED Bulb W3 (Multicolor)}, model \textbf{YLDP005}: 100--240~V AC (50/60~Hz), 0.07~A mains rating, 900~lm, 1700--6500~K tunable white plus RGB, 8~W rated consumption. Commands (on/off, brightness, named colors, warm/cool presets, RGB cycle) were verified against this unit on the same LAN as the Flask host.

\end{document}
